% Presentasi review dua paper: Audebert et al. (2020) dan Kashyap et al. (2026).
% Isi mengikuti ../laporan/bagian/isi.tex; angka disalin dari tabel paper.
% Kompilasi: latexmk -pdf main.tex
\documentclass[aspectratio=169,10pt]{beamer}
\usetheme[progressbar=frametitle,numbering=fraction]{metropolis}
\usepackage[indonesian]{babel}
\usepackage{booktabs,multirow,tikz}
\usetikzlibrary{arrows.meta,positioning,calc}
\graphicspath{{../gambar/}}
\setbeamerfont{caption}{size=\scriptsize}
\newcommand\kecil{\fontsize{8}{9.5}\selectfont}
\definecolor{merah}{HTML}{B03A2E}
\newcommand\awas[1]{\textcolor{merah}{\textbf{#1}}}

\title{Review Paper: Klasifikasi Dokumen Multimodal\\Teks OCR + Citra}
\subtitle{Audebert et al.\ (2020) \textbullet{} Kashyap et al.\ (2026)}
\author{Yulian Anggreaif Rudamaga, Zalzabila Rani, Ramadhani Al-Qadri}
\institute{Analisis \& Pemrosesan Data Multimodal -- Magister Ilmu Komputer, FILKOM UB}
\date{}

\begin{document}
\maketitle

\begin{frame}{Dua paper yang direview}
  \kecil
  \begin{tabular}{p{0.15\linewidth}p{0.38\linewidth}p{0.38\linewidth}}
    \toprule
    & \textbf{Paper 1} & \textbf{Paper 2} \\ \midrule
    Judul & Multimodal Deep Networks for Text and Image-Based Document Classification
          & Multimodal Fusion and Supervised Contrastive Learning for Robust Document Image Classification \\[2pt]
    Penulis & Audebert, Herold, Slimani, Vidal (Quicksign) & Kashyap, Mittal, Singh (Central Univ.\ of Punjab) \\[2pt]
    Venue & ECML PKDD 2019 Workshops, CCIS 1167, Springer 2020 & Int.\ J.\ Intelligent Systems (Wiley), 2026 \\[2pt]
    DOI & 10.1007/978-3-030-43823-4\_35 & 10.1155/int/4241437 \\[2pt]
    Peran & Dasar replikasi Fase 1 & Referensi utama Fase 2 \\ \bottomrule
  \end{tabular}

  \vspace{6pt}
  Aspek review: 9 poin wajib (latar belakang, dataset, modalitas, pipeline, model unimodal,
  fusi, metrik, hasil, future work) + 4 tambahan (kontribusi, detail pelatihan, kelebihan,
  \alert{catatan kritis}).
\end{frame}

% =====================================================================
\section{Paper 1 -- Audebert et al.\ (2020)}

\begin{frame}{Latar belakang, dataset, modalitas}
  \begin{columns}[T]
    \column{0.5\textwidth}
    \textbf{Masalah}
    \begin{itemize}\kecil
      \item Dokumen hasil pindai/foto tidak punya lapisan teks; pengenalan jenis dokumen masih manual.
      \item Visual saja cukup untuk kelas kasar, tidak untuk kelas halus (formulir pajak: tata letak sama, isi beda).
      \item Pertanyaan: apakah teks OCR yang berderau tetap menambah akurasi jika digabung dengan citra, dengan biaya rendah?
    \end{itemize}
    \column{0.5\textwidth}
    \textbf{Dataset}
    \begin{itemize}\kecil
      \item Tobacco3482: 3.482 dok., 10 kelas; 800 latih, sisanya uji; rerata 3 run.
      \item RVL-CDIP: 400 rb dok., 16 kelas; split 320k/40k/40k.
      \item QS-OCR: teks Tesseract kedua dataset, dirilis terbuka.
    \end{itemize}
    \textbf{Modalitas}
    \begin{itemize}\kecil
      \item Citra grayscale halaman + teks OCR dari citra yang sama.
      \item Teks adalah turunan citra; tata letak tidak dipakai.
    \end{itemize}
  \end{columns}
\end{frame}

\begin{frame}{Pipeline}
  % Pipeline Audebert et al. (2020), digambar ulang dari Gambar 1 dan Bagian 3-4 paper.
% Dipakai bersama oleh laporan dan presentasi; butuh TikZ (positioning, arrows.meta).
\resizebox{\linewidth}{!}{%
\begin{tikzpicture}[
  node distance=4mm and 7mm,
  kotak/.style={draw, rounded corners=2pt, align=center, font=\small,
                minimum height=10mm, inner sep=3pt},
  citra/.style={kotak, fill=blue!10},
  teks/.style={kotak, fill=orange!15},
  fusi/.style={kotak, fill=green!15},
  panah/.style={-{Stealth[length=2mm]}, thick}]
  \node[kotak] (in) {Citra dokumen\\(grayscale)};
  % alur citra
  \node[citra, above right=4mm and 9mm of in] (rs) {Resize 384$\times$384\\gray$\to$3 kanal};
  \node[citra, right=of rs] (mb) {MobileNetV2\\(ImageNet, fine-tune)};
  \node[citra, right=of mb] (gap) {GAP 1280-d\\$\to$ FC 128-d};
  % alur teks
  \node[teks, below right=4mm and 9mm of in] (ocr) {Tesseract 4.0\\(LSTM, EN)};
  \node[teks, right=of ocr] (ft) {SpaCy token\\+ FastText (Common Crawl)};
  \node[teks, right=of ft] (cnn) {CNN 1D 4 lapis\\$\to$ 128-d};
  % fusi
  \node[fusi, right=12mm of $(gap)!0.5!(cnn)$] (cat) {Konkatenasi\\256-d};
  \node[fusi, right=of cat] (out) {MLP +\\softmax};
  \draw[panah] (in) |- (rs);  \draw[panah] (rs) -- (mb);  \draw[panah] (mb) -- (gap);
  \draw[panah] (in) |- (ocr); \draw[panah] (ocr) -- (ft); \draw[panah] (ft) -- (cnn);
  \draw[panah] (gap) -| (cat); \draw[panah] (cnn) -| (cat); \draw[panah] (cat) -- (out);
  \node[font=\footnotesize\itshape, below=1mm of out, align=center] {dilatih\\end-to-end};
\end{tikzpicture}}

  \vspace{6pt}
  \begin{itemize}\kecil
    \item Dua cabang menghasilkan vektor 128-d, lalu dikonkatenasi dan diklasifikasi MLP.
    \item Seluruh jaringan dilatih end-to-end (FastText tetap beku).
  \end{itemize}
\end{frame}

\begin{frame}{Model unimodal dan fusi}
  \begin{columns}[T]
    \column{0.5\textwidth}
    \textbf{Citra: MobileNetV2}
    \begin{itemize}\kecil
      \item Ringan; akurasi setara VGG-16, jauh lebih cepat.
      \item 384$\times$384, grayscale diduplikasi 3 kanal, bobot ImageNet.
      \item Augmentasi tidak membantu (84,5 vs.\ 83,9\,\%).
    \end{itemize}
    \textbf{Teks: FastText + CNN 1D}
    \begin{itemize}\kecil
      \item Tesseract 4.0 LSTM; $\sim$26\,\% kata tidak ada di kamus.
      \item FastText (subkata) agar kata salah eja tetap punya vektor.
      \item CNN 1D: 4 lapis, kernel 12, 512 kanal, max-pool waktu; 500 kata.
      \item MLP embedding dokumen kalah (70,8 vs.\ 73,9\,\%).
    \end{itemize}
    \column{0.5\textwidth}
    \textbf{Fusi: tingkat fitur}
    \begin{itemize}\kecil
      \item Teks 128-d $\oplus$ citra (GAP 1280 $\to$ FC 128) $\to$ MLP.
      \item Penjumlahan fitur dicoba, tetapi kalah dari citra saja.
    \end{itemize}
    \textbf{Pelatihan}
    \begin{itemize}\kecil
      \item SGD momentum 0,9; lr 0,01; batch 40.
      \item Epoch: teks 100, citra 200, fusi 200.
      \item TensorFlow 1.12, Titan X.
    \end{itemize}
    \textbf{Metrik}
    \begin{itemize}\kecil
      \item OA, macro-F1, F1/kelas, \alert{oracle}, latensi.
    \end{itemize}
  \end{columns}
\end{frame}

\begin{frame}{Hasil}
  \kecil
  \textbf{Tobacco3482} (OA \%, F1 per kelas)\\[2pt]
  \setlength\tabcolsep{3pt}
  \begin{tabular}{lcccccccccccc}
    \toprule
    Model & OA & F1 & Adv & Email & Form & Letter & Memo & News & Note & Report & Res & Sci \\ \midrule
    Teks & 73,8 & 0,71 & 0,60 & 0,96 & 0,76 & 0,71 & 0,79 & 0,67 & 0,62 & 0,43 & \textbf{0,97} & 0,57 \\
    Citra & 84,5 & 0,82 & \textbf{0,94} & 0,96 & 0,85 & 0,83 & 0,90 & 0,89 & 0,83 & 0,61 & 0,80 & 0,62 \\
    \textbf{Fusi} & \textbf{87,8} & \textbf{0,86} & 0,93 & 0,98 & 0,88 & 0,86 & 0,90 & 0,90 & 0,85 & 0,71 & 0,96 & 0,68 \\
    Oracle & 92,1 & 0,91 & 0,94 & 0,99 & 0,94 & 0,92 & 0,93 & 0,93 & 0,89 & 0,81 & 0,97 & 0,79 \\ \bottomrule
  \end{tabular}

  \vspace{8pt}
  \begin{columns}[T]
    \column{0.48\textwidth}
    \textbf{RVL-CDIP} (OA \%)\\[2pt]
    \begin{tabular}{ccc}
      \toprule Citra & Teks & Fusi \\ \midrule 89,1 & 74,6 & \textbf{90,6} \\ \bottomrule
    \end{tabular}
    \column{0.52\textwidth}
    \begin{itemize}
      \item Fusi +3,3 (Tobacco) dan +1,5 (RVL) di atas citra.
      \item Oracle: batas atas +7,6, jadi masih ada celah 4,3.
      \item Latensi CPU: OCR 910 ms + model 360 ms $\approx$ 1,3 s; OCR jadi hambatan utama.
    \end{itemize}
  \end{columns}
\end{frame}

\begin{frame}{Future work, kelebihan, catatan kritis}
  \begin{columns}[T]
    \column{0.33\textwidth}
    \textbf{Future work}
    \begin{itemize}\kecil
      \item Data dunia nyata (foto, rotasi) + augmentasi.
      \item Segmentasi tata letak (dhSegment).
      \item Koreksi ejaan (SymSpell) + NER.
      \item Embedding tahan derau OCR.
    \end{itemize}
    \column{0.31\textwidth}
    \textbf{Kelebihan}
    \begin{itemize}\kecil
      \item Desain sederhana, tiap pilihan diuji.
      \item Oracle dan F1 per kelas.
      \item Analisis derau OCR kuantitatif.
      \item Data teks terbuka sehingga mudah direplikasi.
    \end{itemize}
    \column{0.36\textwidth}
    \textbf{Catatan kritis}
    \begin{itemize}\kecil
      \item Tanpa simpangan baku dan uji signifikansi.
      \item Cara memilih 800 dokumen latih tidak dijelaskan.
      \item OA teks 73,9 vs.\ 73,8 (Tabel 1 vs.\ 3).
      \item RVL-CDIP hanya OA.
      \item \alert{Satu OCR saja}; tidak ada pembanding late fusion.
    \end{itemize}
  \end{columns}
\end{frame}

% =====================================================================
\section{Paper 2 -- Kashyap et al.\ (2026)}

\begin{frame}{Latar belakang, dataset, modalitas}
  \begin{columns}[T]
    \column{0.5\textwidth}
    \textbf{Masalah}
    \begin{itemize}\kecil
      \item Unimodal tidak tangguh: visual ambigu, teks OCR penuh OOV.
      \item VLM besar (LayoutLMv3, DocFormer, DocLLM) terlalu berat untuk edge/IoMT.
      \item Fusi awal/menengah rentan propagasi galat antarmodal.
      \item Celah: trade-off akurasi--efisiensi jarang dibahas, sehingga diusulkan \alert{dua kerangka} terpisah.
    \end{itemize}
    \column{0.5\textwidth}
    \textbf{Dataset} (semua model dilatih di RVL-CDIP)\\[2pt]
    {\kecil\begin{tabular}{lrrc}
      \toprule Dataset & Total & Uji & Kelas \\ \midrule
      RVL-CDIP & 399.743 & 39.995 & 16 \\
      Tobacco3482 & 3.482 & 3.482 & 10 \\
      RealDoc-16 (primer) & 160 & 160 & 16 \\ \bottomrule
    \end{tabular}}
    \vspace{4pt}
    \textbf{Modalitas}
    \begin{itemize}\kecil
      \item Fusi: citra + teks OCR.
      \item SCSR: teks saja.
    \end{itemize}
  \end{columns}
\end{frame}

\begin{frame}{Pipeline}
  \centering
  \begin{minipage}{0.85\linewidth}% Dua pipeline Kashyap et al. (2026), digambar ulang dari Gambar 1 dan Gambar 3 paper.
% Dipakai bersama oleh laporan dan presentasi; butuh TikZ (positioning, arrows.meta).
\resizebox{\linewidth}{!}{%
\begin{tikzpicture}[
  node distance=4mm and 6mm,
  kotak/.style={draw, rounded corners=2pt, align=center, font=\small,
                minimum height=10mm, inner sep=3pt},
  citra/.style={kotak, fill=blue!10},
  teks/.style={kotak, fill=orange!15},
  fusi/.style={kotak, fill=green!15},
  ret/.style={kotak, fill=violet!12},
  panah/.style={-{Stealth[length=2mm]}, thick}]
  % ---- Pendekatan 1: late fusion
  \node[font=\small\bfseries, anchor=west] (j1) at (0,1.9) {Pendekatan 1 -- multimodal late fusion};
  \node[kotak] (in) at (0.9,0) {Citra\\dokumen};
  \node[citra, above right=1mm and 8mm of in] (pre) {Resize 300$\times$300\\norm.\ [0,1], RGB};
  \node[citra, right=of pre] (eff) {EfficientNetB3\\fine-tune RVL-CDIP};
  \node[citra, right=of eff] (lv) {logit visual\\16 kelas};
  \node[teks, below right=1mm and 8mm of in] (ocr) {OCR komersial\\+ praproses teks};
  \node[teks, right=of ocr] (bert) {BERT\textsubscript{base}\\fine-tune RVL-CDIP};
  \node[teks, right=of bert] (lt) {logit teks\\16 kelas};
  \node[fusi, right=10mm of $(lv)!0.5!(lt)$] (avg) {Rata-rata logit\\$(z_v+z_t)/2$};
  \node[fusi, right=of avg] (y1) {argmax};
  \draw[panah] (in) |- (pre); \draw[panah] (pre) -- (eff); \draw[panah] (eff) -- (lv);
  \draw[panah] (in) |- (ocr); \draw[panah] (ocr) -- (bert); \draw[panah] (bert) -- (lt);
  \draw[panah] (lv) -| (avg); \draw[panah] (lt) -| (avg); \draw[panah] (avg) -- (y1);
  % ---- Pendekatan 2: SCSR
  \node[font=\small\bfseries, anchor=west] (j2) at (0,-2.4) {Pendekatan 2 -- SCSR (teks saja)};
  \node[teks] (ocr2) at (1.2,-3.6) {Teks OCR\\RVL-CDIP train};
  \node[ret, right=of ocr2] (pair) {319.625 pasangan\\anchor--positif};
  \node[ret, right=of pair] (bge) {bge-small-en\\MNRL + Matryoshka};
  \node[ret, right=of bge] (idx) {chunk 512 kar.\\FAISS 1,03 jt $\times$ 384-d};
  \node[ret, right=of idx] (knn) {top-10 kosinus\\voting mayoritas};
  \draw[panah] (ocr2) -- (pair); \draw[panah] (pair) -- (bge);
  \draw[panah] (bge) -- (idx); \draw[panah] (idx) -- (knn);
\end{tikzpicture}}
\end{minipage}
  \begin{itemize}\kecil
    \item Pendekatan 1 mengejar akurasi: dua aliran dilatih terpisah, lalu logit dirata-rata.
    \item Pendekatan 2 mengejar efisiensi: pada dasarnya $k$-NN ($k=10$) di ruang embedding hasil pelatihan kontrastif.
  \end{itemize}
\end{frame}

\begin{frame}{Model unimodal, fusi, metrik}
  \begin{columns}[T]
    \column{0.5\textwidth}
    \textbf{Citra: EfficientNetB3}
    \begin{itemize}\kecil
      \item 300$\times$300, RGB, [0,1]; 1536 $\to$ Dropout $\to$ 16.
      \item 10 epoch, batch 16, Adam $10^{-3}$; $\approx$22,5 jam.
    \end{itemize}
    \textbf{Teks: BERT\textsubscript{base}}
    \begin{itemize}\kecil
      \item \awas{OCR komersial tanpa nama}; praproses: huruf kecil, buang angka/stopword/tanda baca.
      \item 10 epoch, batch 16, Adam $5\!\times\!10^{-5}$, dropout 0,3; $\approx$41 jam.
    \end{itemize}
    \textbf{Teks: SCSR (bge-small-en, 33,3 jt param.)}
    \begin{itemize}\kecil
      \item MNRL + Matryoshka; 4 epoch, AdamW $2\!\times\!10^{-5}$.
      \item Chunk 512 karakter, FAISS 1,03 jt $\times$ 384-d.
    \end{itemize}
    \column{0.5\textwidth}
    \textbf{Fusi: late fusion} \quad $\hat y = \arg\max_c \tfrac12 (z_v + z_t)_c$
    \begin{itemize}\kecil
      \item Rata-rata logit tak berbobot; tidak ada parameter fusi.
    \end{itemize}
    \textbf{Metrik}
    \begin{itemize}\kecil
      \item Acc; P/R/F1 \emph{macro} (RVL, RealDoc) dan \emph{weighted} (Tobacco).
      \item Akurasi per kelas; parameter, latensi, throughput.
      \item UMAP, Grad-CAM, integrated gradients.
    \end{itemize}
  \end{columns}
\end{frame}

\begin{frame}{Hasil}
  \kecil
  \resizebox{\linewidth}{!}{%
  \begin{tabular}{llcccccccc}
    \toprule
    & & \multicolumn{4}{c}{RVL-CDIP (macro)} & \multicolumn{4}{c}{Tobacco3482 (weighted)} \\
    \cmidrule(lr){3-6}\cmidrule(lr){7-10}
    Studi & Mod. & Acc & F1 & Rec & Prec & Acc & F1 & Rec & Prec \\ \midrule
    Audebert (2019) & T / V / T+V & 74,6 / 89,1 / 90,6 & & & & 73,8 / 84,5 / 87,8 & & & \\
    Zingaro (2021) & T+V & 92,7 & & & & 90,5 & & & \\
    Dauphinee (2019) & T+V & 93,03 & & & & -- & & & \\ \midrule
    \multirow{3}{*}{Kashyap: fusi} & T (BERT) & 89,19 & 89,57 & 89,65 & 90,19 & 90,28 & 92,62 & 90,28 & 97,28 \\
    & V (EffNet) & 89,13 & 89,12 & 89,15 & 89,23 & 89,05 & 91,58 & 89,05 & 96,15 \\
    & \textbf{T+V} & \textbf{93,06} & 93,06 & 93,17 & 93,06 & \textbf{92,00} & 94,02 & 92,00 & 98,23 \\
    Kashyap: SCSR & T & 80,76 & 84,00 & 81,00 & 84,00 & 87,89 & 90,00 & 88,00 & 95,00 \\ \bottomrule
  \end{tabular}}

  \vspace{6pt}
  \begin{itemize}
    \item Teks $\approx$ citra; fusi +3,87 (RVL) dan +1,72 (Tobacco) di atas unimodal terbaik.
    \item Latensi: fusi 0,184 s/dok, SCSR 0,157 s/dok (OCR tidak dihitung).
    \item RealDoc-16: hanya 6 dari 16 kelas dilaporkan (mis.\ letter 100/90, memo 72,73/63,60).
  \end{itemize}
\end{frame}

\begin{frame}{Future work dan kelebihan}
  \begin{columns}[T]
    \column{0.5\textwidth}
    \textbf{Future work}
    \begin{itemize}
      \item Dokumen multibahasa dengan OCR yang sangat terdegradasi.
      \item PEFT untuk menekan memori pelatihan.
      \item Pembobotan modalitas adaptif / fusi berbasis keyakinan.
      \item Tambah embedding citra ringan ke indeks SCSR.
    \end{itemize}
    \column{0.5\textwidth}
    \textbf{Kelebihan}
    \begin{itemize}
      \item Jurnal 2026, open access (memenuhi syarat dosen).
      \item Trade-off akurasi--efisiensi eksplisit.
      \item Uji di tiga dataset, termasuk data primer.
      \item P/R/F1 + per kelas untuk tiap unimodal.
      \item Late fusion mudah direplikasi.
    \end{itemize}
  \end{columns}
\end{frame}

\begin{frame}{Catatan kritis}
  \begin{columns}[T]
    \column{0.5\textwidth}
    \begin{itemize}\kecil
      \item \awas{OCR tidak disebut}, sehingga tidak bisa direplikasi.
      \item \awas{Pemetaan kelas Tobacco hilang}: \emph{note} + \emph{report} (13,4\,\%) tidak ada di RVL-CDIP, sehingga tanpa pemetaan akurasi maks.\ $\approx$86,6\,\%, padahal dilaporkan 92\,\%.
      \item Potensi kebocoran: Tobacco3482 dan RVL-CDIP dari arsip yang sama; tidak disebut ada deduplikasi.
      \item Protokol Tobacco berbeda dari Audebert (transfer vs.\ 800 latih).
      \item Hanya rata-rata tak berbobot; tidak ada ablasi fusi dan oracle.
    \end{itemize}
    \column{0.5\textwidth}
    \begin{itemize}\kecil
      \item \awas{Angka tidak konsisten}: macro-F1 SCSR 84 $>$ (P 84 + R 81)/2; macro-recall $\neq$ akurasi pada data seimbang.
      \item RealDoc-16: 160 vs.\ 120 dokumen; 72,73\,\% = 8/11 padahal 10/kelas.
      \item Angka per kelas Audebert di Tabel 5 \awas{tidak ditemukan} di paper asli (yang melaporkan F1, bukan akurasi).
      \item Klaim ``tanpa DL'' untuk SCSR padahal transformer 33 jt; indeks FAISS $\approx$1,6 GB; hemat latensi hanya $\sim$15\,\%.
      \item Tanpa seed, variansi, atau uji signifikansi.
    \end{itemize}
  \end{columns}
\end{frame}

% =====================================================================
\section{Perbandingan dan relevansi}

\begin{frame}{Perbandingan}
  \kecil
  \begin{tabular}{p{0.17\linewidth}p{0.37\linewidth}p{0.37\linewidth}}
    \toprule
    Aspek & Audebert (2020) & Kashyap (2026) \\ \midrule
    OCR & Tesseract 4.0, terbuka & Komersial, anonim \\
    Citra & MobileNetV2 & EfficientNetB3 \\
    Teks & FastText (beku) + CNN 1D & BERT\textsubscript{base}; bge-small (SCSR) \\
    Fusi & Fitur, konkatenasi + MLP & Keputusan, rata-rata logit \\
    Unimodal & Citra $\gg$ teks & Teks $\approx$ citra \\
    Fusi Tobacco & 87,8 (+3,3) & 92,00 (+1,72) \\
    Fusi RVL & 90,6 (+1,5) & 93,06 (+3,87) \\
    Protokol Tobacco & 800 latih & Transfer dari RVL-CDIP \\
    Replikasi & Tinggi & Sedang \\ \bottomrule
  \end{tabular}

  \vspace{6pt}
  \alert{Inti:} lompatan kualitas teks di Kashyap berasal dari OCR \emph{dan} model teks
  yang berubah bersamaan, sehingga keduanya tidak bisa dipisahkan.
\end{frame}

\begin{frame}{Relevansi untuk proyek (Fase 2)}
  \begin{itemize}
    \item Kedua paper tidak memisahkan pengaruh \textbf{OCR} dari pengaruh \textbf{embedding};
          eksperimen faktorial OCR $\times$ embedding pada proyek kami mengisi celah ini.
    \item Late fusion (rata-rata logit) murah, sehingga bisa ditambahkan sebagai pembanding
          konkatenasi Audebert, bersama \alert{oracle} sebagai batas atas.
    \item Praktik yang kami ambil: rerata $\pm$ simpangan baku lintas seed, macro-F1 untuk
          data tidak seimbang, protokol split dan mesin OCR dinyatakan eksplisit, dan tidak
          membandingkan angka dari protokol yang berbeda.
  \end{itemize}
\end{frame}

\begin{frame}{Kesimpulan}
  \begin{itemize}
    \item \textbf{Audebert}: baseline multimodal sederhana dan transparan; kelemahannya satu OCR
          dan tanpa variansi.
    \item \textbf{Kashyap}: komponen modern dan late fusion sederhana dengan akurasi tinggi,
          ditambah alternatif retrieval ringan, tetapi pelaporannya punya beberapa masalah
          (OCR anonim, pemetaan kelas, angka tidak konsisten).
    \item Angka Kashyap dipakai sebagai \emph{acuan arah}, bukan target yang harus dicapai persis.
  \end{itemize}
\end{frame}

\begin{frame}[standout]
  Terima kasih
\end{frame}
\end{document}
